\section{Validation and Benchmark}
\label{sec:benchmark}

% PAGE TARGET: ~8 pages.
% Day 5 (5/28-5/29). Narrative: synthetic phantoms first (clean ground
% truth, sweeps), then real folios (spreadsheet GT, manual GT), then
% failure modes.

\subsection{Synthetic Phantom Characterisation}
\label{sec:benchmark:phantom}

The benchmark begins with a synthetic phantom rather than a real
folio because the real folios offer no ground truth at the level of
the individual detector parameter. On a real folio the period, the
orientation, the wire width and the signal-to-noise ratio are all
unknown and they all vary together, so a single density figure that
disagrees with a hand count tells the reader nothing about which of
the four contributed to the disagreement. A phantom inverts the
relationship. The image is generated by a routine that places a
periodic train of Gaussian wire shadows at a controlled period,
orientation, wire full-width and additive noise level, and then
runs the pipeline of \Cref{sec:multiphi} on the result. The
phantom is a single greyscale frame, so the run drives the
single-image MSI variant of \Cref{sec:multiphi:single}; the common
preprocessing and evaluation framework of
\Cref{sec:multiphi:common} is exercised on every sweep point. Each
of the four parameters can therefore be swept independently while
the others are held fixed, and the recovered value can be compared
against the value that was actually drawn.

Four sweeps were carried out. The period sweep covers true periods
from 10 to 60 pixels at five random seeds per period. The
signal-to-noise sweep covers SNR values from 2 to 50 at a fixed
period of 24 pixels. The orientation sweep covers true line angles
from $-85^{\circ}$ to $+90^{\circ}$ in steps of five degrees. The
wire-width sweep covers Gaussian wire sigmas from one to eight
pixels at a fixed period of 24 pixels. The headline number for each
sweep is the mean absolute error of the recovered quantity
expressed as a percentage of the true value, averaged over the five
random seeds.
\begin{figure}[!htbp]
  \centering
  \begin{subfigure}{0.49\linewidth}
    \includegraphics[width=\linewidth]{phantom_period_sweep}
    \caption{Period sweep.}
    \label{fig:phantom_period}
  \end{subfigure}
  \hfill
  \begin{subfigure}{0.49\linewidth}
    \includegraphics[width=\linewidth]{phantom_snr_sweep}
    \caption{SNR sweep.}
    \label{fig:phantom_snr}
  \end{subfigure}
  \\[0.4em]
  \begin{subfigure}{0.49\linewidth}
    \includegraphics[width=\linewidth]{phantom_angle_sweep}
    \caption{Orientation sweep.}
    \label{fig:phantom_angle}
  \end{subfigure}
  \hfill
  \begin{subfigure}{0.49\linewidth}
    \includegraphics[width=\linewidth]{phantom_sigma_sweep}
    \caption{Wire-width sweep.}
    \label{fig:phantom_sigma}
  \end{subfigure}
  \caption{Synthetic phantom characterisation; mean of five seeds per point.}
  \label{fig:phantom_sweeps}
\end{figure}

The period sweep is the most important of the four because the
laid-line density is the headline output of the pipeline. Across
the range from 10 to 50 pixels per period the mean absolute period
error stays inside $\pm 2.5\,\%$, with no systematic trend across
the range. The single failure point in the sweep is at 60 pixels,
where the recovered period drops to about half of the true value
and the error rises to $-50\,\%$. The cause is a sub-harmonic
alias closely related to the failure mode that affects one of the
real folios in \Cref{sec:benchmark:failures}: at periods longer
than about a quarter of the image dimension, the first-harmonic
peak in the power spectrum is no longer cleanly separated from
the second-harmonic peak of the immediately shorter period, and
the detector locks onto whichever happens to be larger. The
$\pm 2.5\,\%$ figure is the practical accuracy of the period
estimator in the regime that covers every laid density encountered
in handmade paper from the late medieval period onwards.

The signal-to-noise sweep gives the most reassuring result. From
SNR 50 down to SNR 2 the period error is essentially flat at
$+1.6\,\%$, with no degradation visible at any noise level inside
the swept range. The wire-width error is biased high throughout the
range, by between 70 and 75 per cent, but is also flat with respect
to SNR, indicating that the wire-width bias is a property of the
estimator and not of the noise. The SNR invariance of the period
estimate on a single image reflects the fact that the radial FFT
peak of a periodic signal stays well above the broadband noise
floor across a wide noise range. The multi-phi variant of
\Cref{sec:multiphi:multi} extends this margin further on a real
DRP stack by coherent summation across azimuthal views, but even
in single-image mode the SNR margin on the phantom is large enough
that no degradation is visible at any noise level inside the swept
range.

The orientation sweep confirms that the automatic line-direction
detector described in \Cref{sec:first_attempts:simple} recovers the
laid-line angle to within one degree across almost the entire
swept range. The single exception is the $+90^{\circ}$ case, where
the recovered angle is $+93^{\circ}$ instead of $+90^{\circ}$. The
three-degree error at that one point is a wrap-around artefact of
the modular arithmetic that reduces the recovered angle to the
canonical $(-90^{\circ}, 90^{\circ}]$ interval, not a real loss of
accuracy, and it does not affect the period or wire-width estimates
at any swept angle.

The wire-width sweep shows the most parameter-dependent behaviour.
The recovered full-width-at-half-maximum lies within ten per cent
of the true value for sigmas between two and five pixels, which
covers the wire diameters of every folio in the real benchmark
once the optical scale of \Cref{sec:bg:dataset} is applied. Outside
that range the recovered width departs systematically: at sigma
below two pixels the Gaussian-comb model under-resolves the wire
and reports a width that is too low, and at sigma above five
pixels the wire is wide enough that the inter-wire baseline becomes
ambiguous and the fit drifts. The conclusion is that the FWHM
figure reported alongside each real folio in
\Cref{sec:benchmark:perfolio} can be read as a quantitative wire
diameter for any folio whose laid density falls between five and
fifteen lines per centimetre, but should be read more cautiously
for outliers on either side.

The four sweeps together establish that the pipeline is, on
controlled inputs, accurate to two or three per cent in period,
robust to noise at any plausible level, accurate to within one
degree in orientation, and accurate to within ten per cent in wire
width across the relevant range. These are the figures against which
the real-folio results of the next subsection should be read.

\subsection{Per-folio results on nine real folios}
\label{sec:benchmark:perfolio}

The first real-folio test was conducted on nine folios drawn from
four manuscript stocks held at Cambridge University Library and
catalogued by the Manuscripts Lab as part of the Mapping Paper in
Medieval England corpus. Each folio was imaged with the AddME Lab
directional reflectance microscope at a field of view of
approximately 25 cm wide, large enough to cover the central column
of writing on a typical late medieval folio. The pipeline of
\Cref{sec:multiphi} was applied to each folio using the
single-image MSI variant of \Cref{sec:multiphi:single}, since the
imaging modality for this part of the dataset is a single
transmitted-light TX940IR frame per folio rather than a full DRP
stack. No per-folio tuning beyond the optical scale and the polarity
flag was applied, and the resulting period, wire width and
self-contrast $z$-score were recorded for each folio. The split-half
stability metric is not used in this subsection because it requires
multiple azimuthal views and trivially returns zero variance on
single-image input.

The shelfmark Kk.1.5 f5v is used as the reference case for the
remainder of this subsection because it is the folio on which the
pipeline reports the strongest internal evidence of a successful
detection and for which an independent manual count is also
available. The recovered density on this folio is
8.97 lines per centimetre against a manual count of 9.0 lines per
centimetre, an absolute error of $-0.35\,\%$. The per-segment wire
diameter is 0.39 mm at the median across 16 stripes, with a
$95\,\%$ confidence interval of $[0.35, 0.40]$ mm. The
self-consistency contrast is low at $z = 0.72$, which is the one
soft point in the metric bundle for this folio and is discussed
below.
\begin{figure}[!htbp]
  \centering
  \begin{subfigure}{0.48\linewidth}
    \includegraphics[width=\linewidth]{f5v_overlay}
    \caption{Detected laid-line grid overlaid on the grazing image.}
    \label{fig:f5v_overlay}
  \end{subfigure}
  \hfill
  \begin{subfigure}{0.48\linewidth}
    \includegraphics[width=\linewidth]{interval_Kk1_5_f5v}
    \caption{Inter-peak distance histogram.}
    \label{fig:f5v_interval}
  \end{subfigure}
  \\[0.4em]
  \begin{subfigure}{0.48\linewidth}
    \includegraphics[width=\linewidth]{f5v_wire}
    \caption{Per-segment wire FWHM.}
    \label{fig:f5v_wire}
  \end{subfigure}
  \hfill
  \begin{subfigure}{0.48\linewidth}
    \includegraphics[width=\linewidth]{f5v_split}
    \caption{Split-half (trivial on single-image input).}
    \label{fig:f5v_split}
  \end{subfigure}
  \caption{Detection summary, Kk.1.5 f5v.}
  \label{fig:perfolio_f5v}
\end{figure}

The interval-distribution panel of the Kk.1.5 f5v summary is the
clearest visual indicator that the detection has attached to the
real laid-line signal. The empirical inter-peak distances are
concentrated tightly around the locked period, with a coefficient
of variation just below ten per cent and no secondary mode at any
harmonic of the locked period. The wire-width segment plot shows
that the segment-by-segment FWHM is uniform across the page to
within twenty per cent of the median, which is consistent with a
folio that has neither local damage nor a strong curvature
gradient. The split-half panel is, as noted above, a trivial
sanity check on single-image input.

The self-contrast figure of $z = 0.72$ for Kk.1.5 f5v needs
unpacking. The metric measures how much more periodic the detected
direction is than a random direction on the same folio, and a value
near zero usually flags a folio on which the detector found no
real signal. In this case the low $z$ reflects not the absence of
laid lines but the very smooth surface of the parchment-like paper
on which Kk.1.5 f5v was written: random directions on this folio
also score relatively high because the paper texture itself has
modest periodicity at the wire scale. The companion folio
Kk.1.5 f9v, made from the same stock, reports
$z = 11.96$, in line with the median $z$ across the rest of the
benchmark, and the manual ground truth on both folios agrees with
the pipeline to better than one per cent. This case illustrates a
point that recurs later: no single metric is decisive on its own,
and a folio should be judged by the combination of all five.

The recovered density, per-segment median wire FWHM and
self-contrast $z$-score for each of the remaining eight folios
follow the same pattern: density figures within a few lines per
centimetre of either the spreadsheet stock value or, where
available, the per-folio manual count; wire FWHM figures
between 0.3 and 0.7 mm; and self-contrast $z$-scores that vary
widely (from negative values on Ff.4.9 f42r, where the laid-line
signal is weakest, to over eleven on Kk.1.5 f9v). The remainder
of this section unpacks two observations that motivate the next
two subsections: a striking gap between the manual-count error
and the spreadsheet-count error (\Cref{sec:benchmark:manual_gt}),
and a single folio whose density is reported at less than a
third of the spreadsheet value (\Cref{sec:benchmark:failures}).
% TODO (polish): a summary table with one row per folio listing
% all five quantities would make this passage tighter; the prose
% is sufficient for the supervisor-review draft.

\subsection{Comparison with the spreadsheet ground truth}
\label{sec:benchmark:spreadsheet}

The Mapping Paper in Medieval England corpus is accompanied by a
catalogue spreadsheet maintained by the Manuscripts Lab. The
spreadsheet records, for each manuscript stock in the corpus, a
representative laid density in lines per centimetre, measured by a
researcher under raking light on one folio of the stock. This is
the natural external reference for the present pipeline, because
the same spreadsheet is what a codicologist would consult to
identify the stock of a previously unstudied folio. Eight of the
nine folios in the benchmark have a spreadsheet density recorded
for their stock, and the comparison between the pipeline output
and the spreadsheet value is summarised in \Cref{fig:vs_spreadsheet}.
\begin{figure}[!htbp]
  \centering
  \includegraphics[width=0.7\linewidth]{spreadsheet_scatter}
  \caption{Pipeline density vs ground truth on the nine-folio benchmark.}
  \label{fig:vs_spreadsheet}
\end{figure}

The pipeline reports a density lower than the spreadsheet on
every one of the eight folios for which a spreadsheet value
exists, with absolute errors from $-10\,\%$ on Hh.2.12 f190 to
$-71\,\%$ on Ff.4.15 f24r. Read at face value this would suggest
a systematic detector bias.

The next subsection rules this reading out by introducing a second
independent ground truth -- a per-folio manual count of every
benchmark folio -- against which both the pipeline output and the
spreadsheet figure can be compared. The comparison turns the
nominal pipeline-vs-spreadsheet errors of \Cref{fig:vs_spreadsheet}
into a much more informative picture: the pipeline tracks the
per-folio manual count more closely than the spreadsheet does,
which says less about the pipeline being right and more about
manual counting being intrinsically hard.

\subsection{Per-folio manual ground truth and the difficulty of hand-counting}
\label{sec:benchmark:manual_gt}

A second per-folio reference was assembled by marking, by hand,
every laid wire visible inside a small window on every one of the
nine MSI folios. The annotation used a custom-built tool that
loads the folio image, lets the operator click each wire in turn
along a chosen axis, and reports the mean and median inter-wire
spacing together with the corresponding density in lines per
centimetre. Between ten and sixteen wires were marked per folio.
The result, used as ground truth in the remainder of this
subsection, is summarised alongside the spreadsheet figure and
the pipeline output in \Cref{tab:three_way_gt}.

% TODO: \Cref{tab:three_way_gt} -- 9-row table: folio, manual lpc
% (median), pipeline lpc (mean), spreadsheet lpc, |pipeline-manual|,
% |spreadsheet-manual|. Generated by scripts/compare_vs_manual_gt.py.

The annotation exercise made one observation forcefully which the
report wants to record before the numbers are discussed. Hand
counting laid lines on a calibrated digital image is
substantially harder than the workflow assumed by the catalogue
convention. Each wire must be picked out in a low-contrast
greyscale field; the operator's eye fatigues quickly; the choice
of where to start and where to stop a counting window introduces
a one- or two-wire ambiguity per window; and on folios with heavy
ink coverage or surface damage the eye can lose the periodic
structure entirely for a stretch of wires. The annotation
performed for the present benchmark took the form of a careful
single-pass count over a clear region of each folio; the
catalogue spreadsheet maintained by the Manuscripts Lab was
compiled over months by expert codicologists with on-folio access
under raking light. Both are competent hand-counts. They still
disagree with each other by between $4$ and $32$ per cent on
seven of the nine folios, with a median absolute disagreement of
roughly twelve per cent. The reading the report takes from this
is that two careful hand-counts on the same paper can be expected
to differ by ten per cent or so, and that any single hand-counted
density figure should be treated as carrying an intrinsic
inter-counter uncertainty of that order. This is the difficulty
of the underlying task speaking, not the carelessness of either
count.

Against that background, the pipeline-vs-manual comparison is the
most defensible accuracy figure the report can offer. On eight of
the nine folios the pipeline density agrees with the per-folio
manual count to within ten per cent, and on six of the nine the
agreement is within five per cent. The median signed error is
about $-5\,\%$, indicating a small consistent under-count that
the project would like to characterise more tightly with a larger
sample. The ninth folio, Ff.4.15 f24r, is the sub-harmonic alias
failure discussed in \Cref{sec:benchmark:failures}: the pipeline
reports $3.92$ lines per centimetre against a manual count of
$12.08$, an error of $-68\,\%$, and the metric bundle flags this
folio as low-confidence so the failure is not silent.

The pipeline tracks the per-folio manual count more closely than
the catalogue spreadsheet does. On six of the nine folios the
pipeline density is closer to the manual count than the
spreadsheet density is, and on Hh.2.12 f190 the pipeline and the
manual count agree exactly to two decimal places. Across the
benchmark, the pipeline's median absolute error against the
manual count is about five per cent while the spreadsheet's
median absolute error against the same manual count is roughly
twelve per cent: the automated detector reproduces a per-folio
expert hand-count to within roughly half the variability between
two independent expert hand-counts of the same folio.
\begin{figure}[!htbp]
  \centering
  \includegraphics[width=0.85\linewidth]{three_way_comparison}
  \caption{Pipeline, spreadsheet and manual densities per folio on
    the nine-folio benchmark.}
  \label{fig:three_way}
\end{figure}

The reading the report takes from this comparison is the one
that the meeting with Prof Da Rold settled. The spreadsheet
discrepancy of \Cref{sec:benchmark:spreadsheet} is not evidence
of detector bias. It is evidence of how hard hand-counting on
this material is. The pipeline, given the same folio that an
expert would count by hand, produces a number closer to a second
expert hand-count than the first expert hand-count is. The
pipeline is therefore at least as accurate as careful manual
counting on this material, and brings additional properties
that the manual workflow cannot: it is reproducible to a
deterministic numerical figure across runs, it scales to a full
codex in compute time that is small compared with the manual
session that produced any single line of the catalogue
spreadsheet, and it produces the metric bundle of
\Cref{sec:multiphi:evaluation} alongside each density so that a
downstream user can read whether the number should be trusted.

\subsection{Zoomed-in scans and a controlled reproduction}
\label{sec:benchmark:zoomed}

A second group of four datasets complements the nine full-folio
MSI benchmark of \Cref{sec:benchmark:perfolio}. The four datasets
are full DRP acquisitions and are therefore processed by the
multi-phi DRP variant of \Cref{sec:multiphi:multi} rather than by
the single-image variant used for the MSI folios. They are all
acquired at a field-of-view width of 8.65 cm rather than the
25 cm of the MSI benchmark, and they address two questions left
open by the MSI results. The first is whether the pipeline preserves
its accuracy at higher magnification, where finer laid features
become resolvable but the periodic signal occupies a larger
fraction of the image. The second is whether the pipeline reaches
the same headline accuracy on a reference sample whose physical
origin is more tightly characterised than any of the catalogued
manuscript folios. The four datasets are referred to here by their
short serials. Datasets 2, 2b and 4 are zoomed-in scans of two
manuscript folios from the Mapping Paper corpus, while dataset 10
is a sheet of paper made by hand by Professor Da Rold using a
reproduction of a fifteenth-century European mould.

Dataset 10 is the controlled reference. Because the paper was made
by hand under known mould specifications, it is the closest the
benchmark comes to a true ground truth at the level of mould
geometry. The pipeline reports a period of 56.1 pixels on this
dataset, corresponding to a density of 8.31 lines per centimetre
at the optical scale of the acquisition. The five-metric bundle is
healthy: the self-contrast $z$-score is $+4.05$, the
harmonic-fit $R^{2}$ at four harmonics is $0.35$, the split-half
period standard deviation is below 0.001 pixels, and the
per-segment wire FWHM has an interquartile range of $[0.33, 0.43]$
mm around a median of $0.38$ mm. These values are within the
ranges seen on the strongest folios of the main benchmark, so
the pipeline behaves consistently across the change of optical
scale. The density and the wire diameter that the pipeline
reports for dataset 10 are the most defensible per-folio numbers
in this report and are the appropriate primary reference against
which the rest of the benchmark should be read.
\begin{figure}[!htbp]
  \centering
  \begin{subfigure}{0.48\linewidth}
    \includegraphics[width=\linewidth]{darold_overlay}
    \caption{Detected laid-line grid.}
    \label{fig:darold_overlay}
  \end{subfigure}
  \hfill
  \begin{subfigure}{0.48\linewidth}
    \includegraphics[width=\linewidth]{interval_10}
    \caption{Inter-peak distance histogram.}
    \label{fig:darold_interval}
  \end{subfigure}
  \\[0.4em]
  \begin{subfigure}{0.48\linewidth}
    \includegraphics[width=\linewidth]{darold_wire}
    \caption{Per-segment wire FWHM.}
    \label{fig:darold_wire}
  \end{subfigure}
  \hfill
  \begin{subfigure}{0.48\linewidth}
    \includegraphics[width=\linewidth]{darold_selfc}
    \caption{Self-contrast against null directions.}
    \label{fig:darold_selfc}
  \end{subfigure}
  \caption{Detection summary, dataset 10 (Da Rold reproduction).}
  \label{fig:darold_panel}
\end{figure}

The remaining three small-fov datasets illustrate two practical
points about deployment. Datasets 2 and 2b are different crops of
the same source acquisition. The unrestricted full-frame run
(dataset 2) reports a period of 74.5 pixels with a self-contrast
$z$-score of only $+0.36$ and a harmonic-fit $R^{2}$ of $0.002$.
Both metrics are characteristic of a detection that has locked
onto noise rather than a real laid-line signal, in the failure
pattern documented in \Cref{sec:benchmark:failures}. Re-running
the same acquisition with a crop restricted to a quadrant of the
image with light ink coverage (dataset 2b) leaves the recovered
period essentially unchanged but raises the self-contrast
$z$-score to $+3.35$ and the harmonic-fit $R^{2}$ to $0.10$. The
density figure that emerges from dataset 2b is therefore the
trustworthy one for this folio, and the recipe of cropping to a
quieter region before running the pipeline generalises to any
folio whose full-frame self-$z$ is low.

Dataset 4 is a separate zoomed-in folio acquired in the same
session. Here the configuration specifies a period search range
of 0.017 to 0.106 cm, equivalent to 8 to 50 pixels at the
acquisition scale, in order to exclude the long-period band where
the dataset 2 failure occurred. With that constraint the pipeline
reports a period of 40.2 pixels, a density of 11.3 lines per
centimetre, a self-contrast $z$-score of $+2.08$ and a
harmonic-fit $R^{2}$ of $0.095$. The metric bundle is weaker than
on dataset 10 but uniformly stronger than on dataset 2, and the
recovered density is consistent with the densities of the
high-quality folios in the main benchmark. The dataset 4 result
is the clearest demonstration that the per-folio period range, in
combination with a sensible crop, is sufficient to recover a
trustworthy density from a real folio at the zoomed-in scale.

\subsection{Ablation: spectral power versus autocorrelation}
\label{sec:benchmark:ablation}

The automatic line-direction detector of
\Cref{sec:first_attempts:simple} originally scored each candidate
angle by the autocorrelation of the corresponding one-dimensional
projection rather than by its peak spectral power. The change to
spectral power was made part-way through the project after the
autocorrelation criterion was found to favour the broad correlation
peak associated with image-wide low-frequency texture over the
narrow peak at the true laid-line period, and so to bias the
recovered orientation on folios with heavy ink coverage. A full
re-run with both criteria across all real-folio datasets was out
of scope for the present report; what is offered instead, on a
single representative folio, is the observation that the
spectral-power criterion gives a self-contrast $z$-score for the
recovered direction that exceeds the autocorrelation criterion's
score by more than a factor of two, and that on at least one
benchmark folio the autocorrelation criterion locked onto an
orientation offset from the spectral-power result by enough to
change the recovered density at the first significant figure. A
fuller ablation across the whole benchmark is identified in
\Cref{sec:conclusion:future} as a candidate for follow-up work.

\subsection{Failure modes}
\label{sec:benchmark:failures}

Three failure modes were encountered during the benchmark. None of
them is silent: each one degrades at least two of the five
evaluation metrics described in \Cref{sec:multiphi:evaluation}, so
a downstream user looking at the metric bundle can recognise the
failure and treat the corresponding density estimate with the
appropriate caution. This subsection describes the three modes,
identifies the metric pattern that flags each one, and explains
the underlying mechanism where it is known.

The first and most dramatic failure mode is sub-harmonic aliasing.
The folio Ff.4.15 f24r is the clearest example in the benchmark.
Its spreadsheet density is 13.5 lines per centimetre, corresponding
to a true period of approximately 12 pixels at the optical scale
of the acquisition, while the pipeline reports a density of
3.92 lines per centimetre, corresponding to a recovered period of
about 42 pixels. The ratio of recovered to true period is 3.4, so
the detector has locked onto a long-period structure that is not
the laid wires themselves. The metric bundle reflects this
unambiguously: the self-contrast $z$-score is $-0.8$, indicating
that the detected direction is no more periodic than a random
direction on the same image, and the interval distribution is
spread broadly rather than concentrated near a single peak. The
folio is processed by the single-image MSI variant, so the
split-half check is not informative here; the alias is identified
by the other two metrics alone. The same failure mode appears in
controlled form in
the phantom sweep of \Cref{sec:benchmark:phantom} at a true period
of 60 pixels, where the recovered period drops to about half of
the true value. The shared mechanism is that, when the true
fundamental sits close to the bottom of the period search band or
when its harmonic energy exceeds its fundamental energy in the
power spectrum, the detector promotes a harmonic peak above the
true peak. The corresponding mitigation, used for datasets 2b and
4 in \Cref{sec:benchmark:zoomed}, is to narrow the period search
band so that the offending harmonic is excluded from the band
before the peak is read.
\begin{figure}[!htbp]
  \centering
  \includegraphics[width=0.7\linewidth]{ff4_15_overlay}
  \caption{Sub-harmonic alias failure on Ff.4.15 f24r; grid at 42 px vs true ${\approx}12$ px.}
  \label{fig:fail_octave}
\end{figure}

The second failure mode is the polarity ambiguity already discussed
in the context of the simple detector in
\Cref{sec:first_attempts:simple}. The fitted grid can land either
on the dark wire-shadow side or on the bright inter-wire side of
the broadband signal, and the choice depends on the relative
brightness of the grazing image that anchors the polarity. The
multi-phi pipeline resolves this for reflective acquisitions with
the half-period auto-correction described in
\Cref{sec:multiphi:final}, which samples the column-mean at the
grid positions and at the half-period shift and chooses whichever
lands on the darker side. For transmitted-light images, where the
wires appear bright rather than dark, the same logic is inverted
by setting the user-supplied flag \texttt{wire\_is\_darker} to
false. The combination removes the polarity failure as a
practical concern: in the present benchmark, every folio reports
an overlay that lands on the wire shadows as intended. The flag
was added in response to a single early transmitted-light folio on
which the recovered grid was visibly half a period offset; that
folio is now correctly handled and no longer appears as a failure
case.

The third failure mode is low signal. Folios on which heavy ink
coverage, surface curvature, or paper damage destroys most of the
laid-line contrast give a recovered density that is the result of
the detector reaching for whatever is left of the periodic
structure in the image. The folio Ff.4.9 f42r is an example: its
spreadsheet density is 6.0 lines per centimetre and the pipeline
reports 5.27, a $-12\,\%$ error that is in line with the typical
within-stock variation of the benchmark. The self-contrast
$z$-score on this folio is $-1.1$, however, indicating that the
laid-line direction is no more periodic than a random direction.
In a deployment setting, the negative self-$z$ would be the signal
to flag the folio for manual review rather than to trust the
density figure. The lesson is that the pipeline does not silently
return a wrong number on low-signal inputs; it returns a number,
accompanied by a metric bundle that says the number should not be
trusted. The downstream user is responsible for reading the
bundle, and a deployment pipeline can be configured to suppress
density outputs whenever the self-contrast $z$-score falls below
a chosen threshold.
